% =========================================================
%  sections/06-metode.tex
%  Bagian 6 — Metode Pembelajaran
% =========================================================

\section{Metode Pembelajaran}

\subsection{Prinsip Pembelajaran}

Pembelajaran dirancang dengan pendekatan \istilah{Outcome-Based Education}
(OBE) dan berpusat pada mahasiswa (\istilah{student-centered learning}).
Konsekuensi prinsip ini pada pelaksanaan perkuliahan adalah sebagai berikut.

\begin{itemize}
  \item Capaian pembelajaran (Sub-CPMK) ditetapkan lebih dahulu, kemudian
        dipilih materi, aktivitas, dan asesmen yang mengarah pada capaian
        tersebut (prinsip \istilah{constructive alignment}).
  \item Mahasiswa aktif menghasilkan artefak nyata setiap minggu berupa kode
        program, level, atau laporan, bukan hanya menerima penjelasan materi.
  \item Materi teori yang bersifat teknis (khususnya Unity) dipelajari
        mandiri melalui bahan daring, sedangkan waktu tatap muka digunakan
        untuk praktik, diskusi, dan umpan balik.
  \item Setiap aktivitas memiliki instrumen asesmen yang eksplisit, sehingga
        mahasiswa mengetahui kriteria keberhasilan sejak awal.
\end{itemize}

\subsection{Bentuk Pembelajaran dan Alokasi Waktu}

Bobot mata kuliah 3 SKS dilaksanakan dalam 16 minggu. Mengacu pada ketentuan
beban belajar, setiap 1 SKS setara dengan 50 menit tatap muka, 60 menit belajar
terstruktur, dan 60 menit belajar mandiri per minggu.

\begin{table}[htbp]
\centering
\caption{Bentuk pembelajaran dan alokasi waktu}
\label{tab:alokasi}
\begin{tabularx}{\textwidth}{@{}L{3.4cm} Y C{2.6cm} C{3.2cm}@{}}
\toprule
\headrow \textbf{Bentuk} & \textbf{Kegiatan} & \textbf{Alokasi per Minggu} & \textbf{Total 16 Minggu} \\
\midrule
Tatap muka (TM)        & Kuliah interaktif, demonstrasi kode, dan praktikum terbimbing di laboratorium & 150 menit & 2.400 menit (40 jam) \\
Belajar terstruktur (BT) & Tugas praktikum, latihan terarah, tugas baca, dan pengerjaan proyek terjadwal & 180 menit & 2.880 menit (48 jam) \\
Belajar mandiri (BM)   & Eksplorasi dokumentasi resmi, eksperimen kode, dan penyelesaian proyek & 180 menit & 2.880 menit (48 jam) \\
\midrule
\textbf{Total}         & & \textbf{510 menit} & \textbf{8.160 menit (136 jam)} \\
\bottomrule
\end{tabularx}
\end{table}

\subsection{Metode dan Strategi Pembelajaran}

\begin{table}[htbp]
\centering
\caption{Metode pembelajaran, proporsi, dan CPMK yang didukung}
\label{tab:metode}
\begin{tabularx}{\textwidth}{@{}L{3.8cm} Y C{1.7cm} L{2.6cm}@{}}
\toprule
\headrow \textbf{Metode} & \textbf{Penerapan} & \textbf{Proporsi} & \textbf{CPMK} \\
\midrule
Kuliah interaktif dan demonstrasi kode (\istilah{live coding}) &
Dosen memperagakan penulisan kode langkah demi langkah pada proyektor, mahasiswa mengikuti dan memodifikasi kode secara langsung. Dipakai untuk konsep yang belum dikenal (misalnya struktur \istilah{game loop} dan A*). &
20\% & \CPMK{1}--\CPMK{4} \\
\addlinespace
Praktikum terbimbing di laboratorium &
Mahasiswa mengerjakan lembar kerja praktikum pada komputer masing-masing; dosen dan asisten berkeliling memberikan umpan balik per mahasiswa. &
35\% & \CPMK{1}--\CPMK{4} \\
\addlinespace
\istilah{Problem-Based Learning} (PBL) &
Mahasiswa menyelesaikan masalah terbuka, misalnya mencari penyebab objek menembus dinding atau memilih heuristik A* yang tepat, lalu mempresentasikan solusinya. &
15\% & \CPMK{2} \\
\addlinespace
\istilah{Project-Based Learning} (PjBL) &
Mahasiswa mengerjakan dua proyek bertahap: mini proyek AI (Materi 6--9) dan proyek akhir berupa prototipe Unity (Materi 10--15). &
20\% & \CPMK{2}--\CPMK{4} \\
\addlinespace
Diskusi kelompok, \istilah{peer review}, dan \istilah{playtest} &
Kelompok saling menguji hasil pekerjaan menggunakan lembar \istilah{playtest}, memberi umpan balik terstruktur, dan mencatat \istilah{bug} yang ditemukan. &
10\% & \CPMK{3}, \CPMK{4} \\
\bottomrule
\end{tabularx}
\end{table}

\subsection{Peran Dosen dan Mahasiswa}

\begin{table}[htbp]
\centering
\caption{Peran dosen dan mahasiswa dalam pembelajaran}
\label{tab:peran}
\begin{tabularx}{\textwidth}{@{}L{3.6cm} Y@{}}
\toprule
\headrow \textbf{Pihak} & \textbf{Peran} \\
\midrule
Dosen &
Fasilitator dan pemberi umpan balik: menyiapkan lembar kerja praktikum, memperagakan teknik penulisan kode, menilai artefak mahasiswa, serta mengarahkan perbaikan berdasarkan hasil \istilah{playtest}. \\
\addlinespace
Mahasiswa &
Pembelajar aktif: mengerjakan praktikum mingguan, memelihara repositori kode, bekerja dalam kelompok pada proyek, serta melakukan evaluasi dan refleksi terhadap hasil karyanya sendiri. \\
\bottomrule
\end{tabularx}
\end{table}
